% chapters/15/section-03.tex
\section{Memicu Bunyi dari Peristiwa Permainan}
\label{sec:peristiwa-bunyi}

\subsection{Daftar Peristiwa dan Bunyinya}

Pekerjaan menghubungkan bunyi dimulai dari daftar, bukan dari kode. Tuliskan
peristiwa apa saja yang pantas memperoleh tanggapan bunyi, lalu tentukan bunyi
untuk masing-masing. Daftar itu menjadi rencana kerja sekaligus daftar
pemeriksaan Anda pada akhir bab.

\begin{table}[htbp]
\centering
\caption{Daftar peristiwa permainan beserta bunyinya}
\label{tab:peristiwa-bunyi}
{\setlength{\tabcolsep}{4pt}\footnotesize%
\begin{tabularx}{\textwidth}{@{}L{3.4cm} L{2.7cm} C{1.5cm} Y@{}}
\toprule
\headrow \textbf{Peristiwa} & \textbf{Bunyi} & \textbf{Jenis} & \textbf{Pasangan yang terlihat} \\
\midrule
Scene dimulai & musik latar & BGM &
Musik berjalan sejak layar pertama; tidak ada kekosongan sebelum pemain
bergerak. \\
\addlinespace
Pemain melompat & \code{jumpSound} & SFX &
Animasi lompatan pada Bab~XIV, dan posisi pemain yang benar-benar naik. \\
\addlinespace
Pemain mengambil benda & \code{pickupSound} & SFX &
Angka skor pada antarmuka bertambah, dan bendanya menghilang. \\
\addlinespace
Pemain terluka & \code{hurtSound} & SFX &
Jumlah nyawa berkurang, dan ikon nyawa di sudut layar ikut berubah. \\
\addlinespace
Pemain menyentuh tujuan & \code{winSound} & SFX &
Pesan kemenangan, dan musik latar berhenti supaya tidak bersaing. \\
\addlinespace
Permainan berakhir & \code{gameOverSound} & SFX &
Musik latar berhenti, lalu layar akhir muncul. \\
\bottomrule
\end{tabularx}}
\end{table}

Perhatikan kolom terakhir. Ia bukan hiasan. Kolom itu memaksa Anda memikirkan
pasangan yang terlihat bagi setiap bunyi, sebagaimana dibahas pada
Bagian~\ref{sec:peran-bunyi} \cite{unity-leveldesign}. Bunyi yang tidak
memiliki pasangan yang terlihat perlu dipertanyakan: apakah ia benar-benar
diperlukan?

\subsection{Mengapa Bunyi Dipusatkan pada Satu Skrip}

Pada Bab~XIV, setiap benda mengurus tampilannya sendiri. Untuk bunyi, cara itu
justru menyulitkan. Ada dua hal yang harus berlaku bagi seluruh permainan
sekaligus: hanya satu musik yang boleh berjalan, dan perbandingan kekuatan
musik dengan efek suara harus tetap sama di mana pun bunyi itu dipanggil.

Karena itu, bunyi dipusatkan pada satu skrip. Setiap bagian permainan memanggil
skrip itu, dan tidak memainkan bunyinya sendiri-sendiri. Susunan itu
diperlihatkan pada Gambar~\ref{fig:rantai-bunyi}: seluruh lajur bermuara pada
satu titik.

\subsection{Skrip Pemusat Bunyi}

Daftar~\ref{lst:ch15-manajer} memperlihatkan skrip pemusat itu. Perhatikan
empat hal di dalamnya.

\begin{enumerate}
  \item \textbf{Dua sumber bunyi yang terpisah.} \code{musicSource} untuk
        musik latar dan \code{sfxSource} untuk efek suara, sesuai peringatan
        pada Bagian~\ref{sec:peran-bunyi}.
  \item \textbf{Satu salinan yang hidup terus.} \code{Instance} dan
        \code{DontDestroyOnLoad} menjaga agar musik tidak terputus ketika
        berpindah \istilah{scene}, dan salinan yang berlebih dibuang.
  \item \textbf{\code{PlayOneShot} untuk efek suara}, sesuai pembahasan pada
        Bagian~\ref{sec:audioclip-audiosource}.
  \item \textbf{Penggeser nada} yang dijelaskan pada
        Bagian~\ref{sec:waktu-volume}.
\end{enumerate}

\kodecsharp{unity/ch15_manajer_bunyi.cs}{Pemusat bunyi permainan: musik latar dan efek suara}{lst:ch15-manajer}

\subsection{Menghubungkan Peristiwa dengan Bunyinya}

Skrip kedua menghubungkan peristiwa permainan dengan bunyi yang telah
disediakan pemusatnya. Ia ditempatkan pada benda pemain, sebab di sanalah
sebagian besar peristiwa terjadi.

\kodecsharp{unity/ch15_pemicu_bunyi.cs}{Menghubungkan peristiwa permainan dengan bunyinya}{lst:ch15-pemicu}

Bunyi lompat dipanggil oleh skrip kendali pemain, bukan dibaca sendiri oleh
skrip ini. Alasannya sama dengan alasan pada Bab~XIII: masukan pemain dibaca di
satu tempat saja. Tambahkan dua baris berikut pada \code{ch13_kontrol_pemain.cs}
buatan Anda.

\begin{kodekeluaran}
public PemicuBunyi pemicuBunyi;          // tambahkan pada bagian peubah

// Di dalam Update(), pada pemeriksaan tombol lompat:
if (Input.GetButtonDown("Jump"))
{
    jumpRequested = true;
    if (pemicuBunyi != null) pemicuBunyi.Lompat();
}
\end{kodekeluaran}

\begin{catatabox}
Baris \code{if (pemicuBunyi != null)} bukan kelebihan. Rujukan yang belum
dihubungkan pada \istilah{Inspector} bernilai kosong, dan memanggilnya akan
menghentikan permainan dengan kesalahan. Kebiasaan memeriksa sebelum memanggil
itu dipakai berulang-ulang pada skrip bab ini, dan itulah yang membuat
kesalahan pemasangan menjadi kesalahan yang mudah ditemukan, bukan
kegagalan yang membingungkan.
\end{catatabox}
